\documentclass[12pt]{article}

%%% Essential packages
\usepackage[utf8]{inputenc}
\usepackage[T1]{fontenc}
\usepackage{amsmath,amsfonts,amssymb}
\usepackage{graphicx}
\usepackage{booktabs}
\usepackage{float}
\usepackage{setspace}
\usepackage[labelfont=bf]{caption,subcaption}
\usepackage[table,dvipsnames]{xcolor}
\usepackage[colorlinks=true,linkcolor=blue,urlcolor=blue,citecolor=blue]{hyperref}
\usepackage{natbib}
\usepackage[margin=1in]{geometry}

%%% Common math shortcuts
\newcommand{\prob}{\mathbb{P}}
\newcommand{\expec}{\mathbb{E}}
\newcommand{\var}{\mathrm{var}}
\newcommand{\cov}{\mathrm{cov}}
\newcommand{\corr}{\mathrm{corr}}
\newcommand{\E}{\mathbb{E}}
\DeclareMathOperator*{\argmin}{arg\,min}
\DeclareMathOperator*{\argmax}{arg\,max}
\DeclareMathOperator{\Var}{Var}
\DeclareMathOperator{\Cov}{Cov}

%%% Code-style inline text
\def\code#1{\texttt{#1}}

%%% Paper-specific commands
%% Add your custom \newcommand definitions here.
%% Examples:
%%   \newcommand\SR{\mathrm{SR}}   % Sharpe ratio
%%   \newcommand\bmR{\bm R}        % bold return vector


%%%%%%%%%%%%%%%%%%%%%%%%%%%%%%%%%%%%%%%%%%%%%%%%%%%%%%%%
%%% Title block
%%%%%%%%%%%%%%%%%%%%%%%%%%%%%%%%%%%%%%%%%%%%%%%%%%%%%%%%
\title{[REMOVE] Your Report Title Here}
\author{[REMOVE] Student Name \\ [REMOVE] Student ID \\ FINS2026 --- Fintech}
\date{\today}

\begin{document}
\maketitle

%%%%%%%%%%%%%%%%%%%%%%%%%%%%%%%%%%%%%%%%%%%%%%%%%%%%%%%%
%%% Abstract
%%%%%%%%%%%%%%%%%%%%%%%%%%%%%%%%%%%%%%%%%%%%%%%%%%%%%%%%
\begin{abstract}
\noindent
[REMOVE] Write a concise summary of your report here. State the research
question, the data or methodology used, and the main findings. Aim for
100--200 words.
\end{abstract}

\thispagestyle{empty}
\pagebreak
\setcounter{page}{1}
\onehalfspacing


%%%%%%%%%%%%%%%%%%%%%%%%%%%%
%%% Sections in the main body
%%%%%%%%%%%%%%%%%%%%%%%%%%%%

%%%%%%%%%%%%%%%%%%%%%%%%%%%%
% INTRODUCTION SECTION
%%%%%%%%%%%%%%%%%%%%%%%%%%%%

%% BEGIN introduction
\section{Introduction}\label{sec:intro}

[REMOVE] Open with the broad topic and why it matters. What is the
research question or practical problem you are addressing?

[REMOVE] Briefly summarise how prior work has approached this topic.
Cite key references, for example \citet{FamaFrench_1993}.

[REMOVE] State how your analysis contributes or what angle you take.

[REMOVE] Outline the rest of the report: Section~\ref{sec:data}
describes the data, Section~\ref{sec:methodology} presents the
methodology, Section~\ref{sec:results} reports the results, and
Section~\ref{sec:conclusion} concludes.

%% END introduction


%%%%%%%%%%%%%%%%%%%%%%%%%%%%
% DATA SECTION
%%%%%%%%%%%%%%%%%%%%%%%%%%%%

%% BEGIN data
\section{Data}\label{sec:data}

[REMOVE] Describe the data sources, sample period, and any filters or
cleaning steps. Mention the key variables and how they are constructed.
Table~\ref{tab:descriptive} reports summary statistics.

%%% --- TABLE: Descriptive Statistics --- %%%
\begin{table}[!ht]
\caption{[REMOVE] Descriptive Statistics.}
{\footnotesize
[REMOVE] This table reports descriptive statistics for the main
variables. $N$ is the number of observations. SD is the standard
deviation. P5 and P95 denote the 5th and 95th percentiles.}
\vspace{2mm}
\begin{center}
\label{tab:descriptive}
\scalebox{0.85}{%
\begin{tabular}{l rrrrrrr}
\toprule
Variable & $N$ & Mean & Median & SD & P5 & P95 \\
\midrule
[REMOVE] Var 1 & 10{,}000 & 0.05 & 0.03 & 0.12 & $-$0.18 & 0.28 \\
\addlinespace
[REMOVE] Var 2 & 10{,}000 & 1.23 & 0.98 & 0.87 & 0.04 & 2.91 \\
\addlinespace
[REMOVE] Var 3 & 9{,}500 & 0.71 & 0.54 & 0.62 & 0.09 & 1.78 \\
\bottomrule
\end{tabular}
}
\end{center}
\end{table}

%% END data


%%%%%%%%%%%%%%%%%%%%%%%%%%%%
% METHODOLOGY SECTION
%%%%%%%%%%%%%%%%%%%%%%%%%%%%

%% BEGIN methodology
\section{Methodology}\label{sec:methodology}

[REMOVE] Explain the analytical framework or model you use. Include
at least one key equation with proper punctuation:
\begin{equation}\label{eq:model}
  r_{i,t} - r_{f,t}
    = \alpha_i + \beta_i (r_{m,t} - r_{f,t}) + \varepsilon_{i,t},
\end{equation}
where $r_{i,t}$ is the return on asset~$i$, $r_{f,t}$ is the risk-free
rate, and $r_{m,t}$ is the market return. [REMOVE] Describe any
additional steps such as portfolio construction, event windows, or
robustness checks.

%% END methodology


%%%%%%%%%%%%%%%%%%%%%%%%%%%%
% RESULTS SECTION
%%%%%%%%%%%%%%%%%%%%%%%%%%%%

%% BEGIN results
\section{Results}\label{sec:results}

[REMOVE] Present your findings. Reference tables and figures as you
discuss the results. Table~\ref{tab:panel_regression} reports regression
estimates and Figure~\ref{fig:example} shows the main plot.

%%% --- TABLE: Panel Regression --- %%%
\begin{table}[!ht]
\caption{[REMOVE] Regression Results.}
{\footnotesize
[REMOVE] This table reports OLS/panel regression estimates. The
dependent variable is the excess return (\%). $t$-statistics are in
parentheses. $^{***}$, $^{**}$, and $^{*}$ denote significance at the
1\%, 5\%, and 10\% levels. Fixed effects as indicated.}
\vspace{2mm}
\begin{center}
\label{tab:panel_regression}
\scalebox{0.80}{%
\begin{tabular}{l ccc}
\toprule
 & (1) & (2) & (3) \\
\midrule
[REMOVE] Variable A & 0.54$^{***}$ & 0.48$^{***}$ & 0.41$^{**}$ \\
 & (3.41) & (3.12) & (2.67) \\
\addlinespace
[REMOVE] Variable B & & $-$0.22$^{**}$ & $-$0.19$^{*}$ \\
 & & ($-$2.14) & ($-$1.89) \\
\addlinespace
[REMOVE] Variable C & & & 0.13 \\
 & & & (1.34) \\
\midrule
Time FE & No & Yes & Yes \\
Firm FE & No & No & Yes \\
$N$ & 10{,}000 & 10{,}000 & 10{,}000 \\
$R^2$ & 0.03 & 0.08 & 0.15 \\
\bottomrule
\end{tabular}
}
\end{center}
\end{table}

%%% --- TABLE: Comparison / Summary --- %%%
\begin{table}[!ht]
\caption{[REMOVE] Comparison of Portfolio Returns (\%).}
{\footnotesize
[REMOVE] This table compares mean monthly returns across portfolios
formed on the signal variable. $t$-statistics (Newey--West) are in
parentheses. Sample period: [REMOVE] 2010--2024.}
\vspace{2mm}
\begin{center}
\label{tab:comparison}
\scalebox{0.80}{%
\begin{tabular}{l cccc}
\toprule
 & Low & Mid & High & High$-$Low \\
\midrule
\multicolumn{5}{l}{\textbf{Panel A:} Equal-Weighted} \\
\midrule
Mean return & 0.42 & 0.68 & 1.14 & 0.72 \\
 & (1.89) & (3.12) & (4.56) & (2.78) \\
\addlinespace
Alpha (FF3) & 0.11 & 0.31 & 0.78 & 0.67 \\
 & (0.56) & (1.67) & (3.34) & (2.45) \\
\midrule
\multicolumn{5}{l}{\textbf{Panel B:} Value-Weighted} \\
\midrule
Mean return & 0.38 & 0.54 & 0.89 & 0.51 \\
 & (1.67) & (2.56) & (3.78) & (2.12) \\
\addlinespace
Alpha (FF3) & 0.08 & 0.21 & 0.56 & 0.48 \\
 & (0.42) & (1.12) & (2.56) & (1.89) \\
\bottomrule
\end{tabular}
}
\end{center}
\end{table}

%%% --- FIGURE: Example --- %%%
\begin{figure}[h!]
\begin{center}
\fbox{\rule{0pt}{4cm}\rule{0.95\textwidth}{0pt}}
\end{center}
\vspace{-4mm}
\caption{[REMOVE] Title of your figure.}
{\footnotesize
[REMOVE] Describe what the figure shows. State the sample period, any
filters, and how to read the axes.}
\label{fig:example}
\end{figure}

%% END results


%%%%%%%%%%%%%%%%%%%%%%%%%%%%
% CONCLUSION SECTION
%%%%%%%%%%%%%%%%%%%%%%%%%%%%

%% BEGIN conclusion
\section{Conclusion}\label{sec:conclusion}

[REMOVE] Summarise the main findings in two or three sentences.
Discuss any limitations of your analysis and possible extensions for
future work.

%% END conclusion


%%%%%%%%%%%%%%%%%%%%%%%%%%%%
% BIBLIOGRAPHY
%%%%%%%%%%%%%%%%%%%%%%%%%%%%
\clearpage
\bibliographystyle{plainnat}
\bibliography{references}


%%%%%%%%%%%%%%%%%%%%%%%%%%%%
% APPENDIX
%%%%%%%%%%%%%%%%%%%%%%%%%%%%
\clearpage
\appendix

\begin{center}
{\LARGE \textbf{Appendix}}
\end{center}

\renewcommand{\theequation}{A.\arabic{equation}}%
\renewcommand{\thefigure}{A.\arabic{figure}} \setcounter{figure}{0}
\renewcommand{\thetable}{A.\arabic{table}} \setcounter{table}{0}

%% BEGIN appendix-a
\section{Additional Results}\label{app:a}

[REMOVE] Include any supplementary tables, figures, or robustness checks
here that support the main text but are not essential to the core
argument.

%% END appendix-a

\end{document}
